\documentclass[reqno,11pt]{amsart}
\usepackage{math327-lecture}
\usepackage{needspace}
\newcommand{\Q}{\mathbb{Q}}
\newcommand{\F}{\mathbb{F}}
\lectureinfo{30}{Nov. 27, 2026}{Adjoining Roots and Integral Domains}

% Typeset from handwritten Lecture 32, Nov. 28, 2025, pp. 1--10.
\begin{document}
\raggedbottom
\makelecturetitle



\section{Quotients by a monic polynomial}

Let $R$ be a ring and $R[x]$ the polynomial ring over $R$.
Let $f(x)\in R[x]$ be monic of degree $n$:
\[
  f(x)=x^n+a_{n-1}x^{n-1}+\cdots+a_1x+a_0 .
\]
We consider the quotient ring $R[x]/(f)$, where
\[
  (f)=\{gf: g\in R[x]\}
\]
is the principal ideal generated by $f$. Since $f$ is monic, every
nonzero $h\in(f)$ satisfies $\deg h\geq\deg f$.

Last time we showed that $R[x]/(f)$ has a ``basis''
\[
  1,\ x,\ x^2,\ \dots,\ x^{n-1}.
\]
That is, every class $g(x)+(f)$ in $R[x]/(f)$ has a representative of
the form
\[
  g=b_0+b_1x+\cdots+b_{n-1}x^{n-1},\qquad b_i\in R,
\]
and moreover the coefficients $b_i$ are unique.

In terms of these representatives, the ring operations are as follows.
\begin{itemize}
  \item Addition is done componentwise.
  \item Multiplication is done as for polynomials, for instance
  \[
    (b_0+b_1x+b_2x^2+b_3x^3)\cdot x^5
    =b_0x^5+b_1x^6+b_2x^7+b_3x^8,
  \]
  after which high powers of $x$ are reduced using the relation
  coming from $f$: in $R[x]/(f)$,
  \[
    x^n=-a_{n-1}x^{n-1}-\cdots-a_1x-a_0 .
  \]
\end{itemize}

\section{Adjoining a cube root of \texorpdfstring{$2$}{2}}

\begin{example}[The ring {$\Z[\gamma]$}]
Put
\[
  \Z[\gamma]\coloneqq\Z[x]/(x^3-2),
\]
where $\gamma$ denotes the class of $x$, so $\gamma^3=2$.
The ring $\Z[x]/(x^3-2)$ has the basis $1,x,x^2$: its elements have
unique representatives
\[
  b_0+b_1x+b_2x^2,\qquad b_i\in\Z.
\]
For example, using $x^3=2$,
\begin{align*}
  (3+5x)(7x+8x^2)
  &=21x+24x^2+35x^2+40x^3\\
  &=21x+59x^2+40x^3\\
  &=80+21x+59x^2 .
\end{align*}
In the notation of $\gamma$, this reads
$(3+5\gamma)(7\gamma+8\gamma^2)=80+21\gamma+59\gamma^2$.
Similarly, since $x^3=2$, we get $x^6=4$, so $x^7=4x$ and
$x^8=4x^2$. Hence
\[
  (3+5x)x^7=3x^7+5x^8=12x+20x^2 .
\]
\end{example}

\begin{example}[The field $\Q(\gamma)$]
Now work over $\Q$ instead:
\[
  \Q[\gamma]\coloneqq\Q[x]/(x^3-2).
\]
Its elements are uniquely of the form
\[
  b_0\cdot1+b_1x+b_2x^2,\qquad b_i\in\Q .
\]
Thus $\Q[\gamma]$ is a $3$-dimensional $\Q$-vector space. But it also
has a multiplication, which is not visible from linear algebra alone.
\end{example}

\begin{proposition}\label{prop:Q-gamma-field}
The ring $\Q[\gamma]=\Q[x]/(x^3-2)$ is a field.
\end{proposition}
\begin{proof}
The ring $\Q[x]$ is a \emph{principal ideal domain}: every ideal of
$\Q[x]$ has the form $(g)$ for some $g\in\Q[x]$.

Ideals $I$ of $\Q[x]/(x^3-2)$ correspond to ideals of $\Q[x]$ that
contain $x^3-2$. Such an ideal has the form $(g)$ with
\[
  (g)\supseteq(x^3-2),
\]
which happens if and only if $g\mid x^3-2$ in $\Q[x]$.
Now $x^3-2$ is irreducible over $\Q$: a factorization would require a
linear factor, hence a rational root, and $2$ has no rational cube
root. So $g$ is either a nonzero constant or a constant multiple of
$x^3-2$. Correspondingly, $(g)=\Q[x]$ or $(g)=(x^3-2)$, and the only
ideals of $\Q[x]/(x^3-2)$ are $(0)$ and the whole ring. A commutative
ring with exactly these two ideals is a field.
\end{proof}
% NOTE: The notes give only the correspondence of ideals and the condition
% g | x^3-2; the irreducibility of x^3-2 and the final step were filled in.

Since $\Q[\gamma]$ is also a field, we sometimes write $\Q(\gamma)$ for
it. We picture it as an extension of degree $3$ over $\Q$:
\[
  \begin{array}{c}
    \Q(\gamma)\\
    \big|\,{\scriptstyle 3}\\
    \Q
  \end{array}
\]

\section{Adjoining square roots over finite fields}

Write $\F_5=\Z/5\Z$.

\begin{example}[The field {$\F_5[\sqrt3]$}]
Consider
\[
  \F_5[x]/(x^2-3)\eqqcolon\F_5[\sqrt3].
\]
Its elements are uniquely of the form $a+bx$ with $a,b\in\F_5$, and
$x$ plays the role of ``$\sqrt3$'' in $\F_5$. Note that $3$ has no
square root in $\F_5$ itself: the squares are
\[
  0^2\equiv0,\quad 1^2\equiv1,\quad 2^2\equiv4,\quad
  3^2\equiv4,\quad 4^2\equiv1 \pmod 5.
\]
So $\F_5[\sqrt3]$ is a ring with $25$ elements, and moreover it is
again a field.

To see this, we show that each nonzero $a+b\sqrt3$ has an inverse in
$\F_5[\sqrt3]$. (Equivalently, the only ideals of $\F_5[\sqrt3]$ are
$(0)$ and $(1)$.) We compute
\[
  (a+b\sqrt3)(a-b\sqrt3)=a^2-3b^2 \in\F_5 .
\]
Here $a^2-3b^2=0$ if and only if $a=b=0$: if $b\neq0$, then
$(a/b)^2=3$, which is impossible by the list of squares above, and
if $b=0$, then $a^2=0$ forces $a=0$. Therefore, for
$a+b\sqrt3\neq0$,
\[
  (a+b\sqrt3)^{-1}=\frac{1}{a^2-3b^2}\,(a-b\sqrt3).
\]
Hence $\F_5[\sqrt3]$ is a \emph{finite} field, with $25$ elements.
\end{example}

\begin{remark}
In general, every finite field has size $p^a$ for some prime $p$ and
positive integer $a$, and is of the form
\[
  \F_p[x]/(f),\qquad f \text{ irreducible of degree } a .
\]
\end{remark}

\begin{example}[A quotient that is not a field]
Consider $\F_{11}[x]/(x^2-3)$. Since $5^2=25\equiv3\pmod{11}$,
\[
  x^2-3=(x-5)(x+5)\quad\text{in }\F_{11}[x].
\]
The ring $\F_{11}[x]/(x^2-3)$ has $\F_{11}$-basis $1,x$, so its
elements are uniquely $a_0+a_1x$ with $a_i\in\F_{11}$.
But it is \emph{not} a field: in the quotient,
\[
  (x-5)(x+5)=x^2-3=0,
\]
while $x-5\neq0$ and $x+5\neq0$ in $\F_{11}[x]/(x^2-3)$.
A field cannot contain two nonzero elements with product zero.
\end{example}

\section{Toward factorization: integral domains}

Recall the main facts about factorization over $\Z$:
\begin{itemize}
  \item the ideals of $\Z$ are exactly the $(n)$;
  \item any two integers $a,b$ have a greatest common divisor
  $\gcd(a,b)$;
  \item if $p$ is a prime number and $p\mid ab$, then $p\mid a$ or
  $p\mid b$;
  \item every $n$ has a prime factorization $n=p_1\cdots p_r$.
\end{itemize}
Similar facts hold for $\Q[x]$. Before discussing the notions of prime
and factorization in a general ring, we need some terminology.

\begin{define}[Integral domain]
A ring $R$ is called an \emph{integral domain} if whenever $a,b\in R$
satisfy $a\cdot b=0$, either $a=0$ or $b=0$.
\end{define}

\begin{proposition}[Cancellation law]\label{prop:cancellation}
In an integral domain, if $ab=ac$ with $a\neq0$, then $b=c$.
\end{proposition}
\begin{proof}
If $ab=ac$, then $a(b-c)=0$. Since $a\neq0$ and $R$ is an integral
domain, $b-c=0$, that is, $b=c$.
\end{proof}

\begin{warningbox}
The cancellation law does not come from dividing by $a$:
in an integral domain we may not be able to divide.
\end{warningbox}

\Needspace{6\baselineskip}
We will compare notions for elements with the corresponding notions
for ideals.

\begin{center}
\begin{tabular}{p{0.45\linewidth}|p{0.4\linewidth}}
  \textbf{Elements} & \textbf{Ideals}\\ \hline
  \rule{0pt}{2.6ex}%
  $u$ is a \emph{unit}, i.e.\ $u\in R^\times$ ($u$ has an inverse):
  there exists $v$ with $u\cdot v=1$.
  & $(u)=(1)=R$.
\end{tabular}
\end{center}

\end{document}
